% Generated draft: requires mathematical and source review.
\documentclass[11pt]{article}
\usepackage[margin=25mm]{geometry}
\usepackage{fontspec}
\setmainfont{Latin Modern Roman}
\usepackage{amsmath,amssymb,mathtools,unicode-math}
\setmathfont{Latin Modern Math}
\usepackage{enumitem,needspace,xurl,hyperref}
\setlength{\parindent}{0pt}
\setlength{\parskip}{5pt}
\setlength{\emergencystretch}{3em}
\newcommand{\N}{\mathbb{N}}
\newcommand{\Z}{\mathbb{Z}}
\newcommand{\Q}{\mathbb{Q}}
\newcommand{\R}{\mathbb{R}}
\newcommand{\C}{\mathbb{C}}
\newcommand{\F}{\mathbb{F}}
\newcommand{\A}{\mathbb{A}}
\newcommand{\PP}{\mathbb{P}}
\newcommand{\E}{\mathbb{E}}
\newcommand{\eps}{\varepsilon}
\newcommand{\dd}{\,\mathrm{d}}
\begin{document}
\section*{UnsolvedMath ID: 800916}
\textbf{Source problem code:} catalogue entry
\subsection{Definitions and mathematical statement}
A \emph{threefold} is an algebraic variety of dimension \(3\). Gromov--Witten (GW) theory concerns invariants defined using moduli spaces of stable maps, while Donaldson--Thomas (DT) theory concerns invariants defined using moduli spaces of sheaves; in the usual curve-counting formulation, the latter include ideal sheaves of subschemes of dimension at most one. These descriptions presuppose geometric hypotheses and choices under which the relevant invariants are defined.

\textbf{Requested statement (Pandharipande, 2010, Q4).} Prove the Gromov--Witten/Donaldson--Thomas correspondence for threefolds.

\textbf{Unresolved formulation.} The supplied record does not specify the class of threefolds, the base field, smoothness or projectivity assumptions, or the precise versions of the invariants to be compared. It also does not specify curve classes, insertions, treatment of degree-zero contributions, generating-series normalization, or change of variables. Consequently, it does not determine a unique fully quantified identity of invariants or generating functions. These choices must be recovered from the cited source before a precise correspondence can be stated; no particular choice is imposed here.

\textbf{Supplied qualifications.} The source record states that the toric cases of Q4 are known, citing the reference label \texttt{[moop]}. It further states that algebraic cobordism results, cited as \texttt{[levp]}, suggest the possibility of reducing to the toric case by degeneration methods. This is a proposed route, not a supplied assertion that such a reduction has been established. The registry labels the problem \emph{open}; this is supplied curation rather than an independently verified assessment of its present status.
\subsection{Short English statement}
Pandharipande asks for a proof of the correspondence between Gromov--Witten and Donaldson--Thomas curve-counting theories for threefolds. The supplied text reports known toric cases and suggests a possible reduction using algebraic cobordism and degeneration. The geometric hypotheses and exact correspondence are unspecified, so the record alone does not determine a fully precise theorem.
\subsection{Sources}
\begin{itemize}
\item Local UnsolvedMath registry, permanent ID \(800916\), exact source problem code \texttt{catalogue entry}; title: \emph{AGNES Open Problems in Algebraic Geometry --- Q4}. The supplied record attributes the problem to Pandharipande, proposed in 2010, and labels it open.
\item Supplied source list: Pandharipande, \emph{Open Problems in Enumerative Algebraic Geometry (AGNES)} (2010), source item Q4 (\#3), \url{https://arxiv.org/abs/1004.3259}. The provenance records extraction by \texttt{arxiv\_latex} and access on 2026-10-16. These are supplied provenance details; the source has not been independently consulted for this draft.
\item Source reference label \texttt{[moop]} accompanies the supplied claim that the toric cases of Q4 are known. No bibliographic entry or separate link for this label is supplied.
\item Source reference label \texttt{[levp]} accompanies the supplied suggestion that algebraic cobordism results may enable reduction to the toric case by degeneration methods. No bibliographic entry or separate link for this label is supplied.
\end{itemize}
\end{document}
